\section{Numerical controls}
\label{sec:numerical-controls}
\subsection{Euler admissibility}
Oscillation damping does not guarantee positive density and pressure. The
implementation therefore includes a separate control for Euler states. For
an element mean $\overline{\boldsymbol u}_K$, it replaces a candidate polynomial
by
\begin{equation}
 \boldsymbol u_h^{\rm lim}
 =\overline{\boldsymbol u}_K+	heta_K
          (\boldsymbol u_h-\overline{\boldsymbol u}_K),
 \qquad 0\leq\theta_K\leq1.
\end{equation}
A single factor is used for all conservative components. Integration shows
that the component means are unchanged. The factor is restricted first for
density, then for pressure by bisection along the segment toward the mean.
Admissibility is checked at the volume and face points needed by the evolution
and filtering operators; it is not a statement about every point of the
continuous polynomial.

A nonphysical mean cannot be repaired by scaling toward itself. In that case
the candidate step is rejected and retried with a smaller timestep, subject to
a retry limit. Comparisons between Euler methods must apply the same
admissibility control and record its interventions so that its influence is
not attributed to the oscillation filter.

\MissingWork{POSITIVITY-SOURCE}{Review the proposed positivity-preserving DG
reference before adding its citation and relating this implementation to the
published assumptions. The scaling formula above describes the implemented
control and does not claim a positivity theorem for the complete scheme.}

\subsection{Experimental choices still to be fixed}
The formulation distinguishes approximation order, spatial resolution, CFL
number, filtering cadence, detector threshold, and positivity treatment.
These must be specified when comparing methods. In particular, comparing
post-step OFDG with stage-wise OEDG changes both the sensor and the placement
of filtering. Such a comparison cannot isolate the sensor's effect without
an additional controlled comparison of cadences.

The final selection of problems, reference solutions, error measures, and
resolution studies is deferred to the Results chapter. Existing exploratory
calculations and their reference-solver workflow remain separate from this
methodological account.
